\documentclass[11pt]{article}
\usepackage[margin=1in]{geometry}
\usepackage{amsmath,amssymb,amsthm}
\newtheorem{theorem}{Theorem}
\title{Exact Fourier correction and its concentration cost}
\author{Joshua Mark}
\date{20 November 2026; research draft, priority unestablished}
\begin{document}\maketitle
Let $P$ be any probability law on the $d$-dimensional sign cube,
$M=2^d$, and $\chi_S(x)=\prod_{i\in S}x_i$. Write
\[
 m_S=\sum_xP(x)\chi_S(x),\qquad
 f(x)=\sum_{1\le|S|\le2}m_S\chi_S(x),\quad
 B=\max_x f(x),\quad L=\sum_{1\le|S|\le2}|m_S|.
\]
All quantities are computable from a finite law. In particular
$0\le B\le L$. Let $C(P)=\sum_xP(x)^2$.

\begin{theorem}[Nonnegative exact correction]
The law
\[
 Q(x)=\frac{P(x)+(B-f(x))/M}{1+B}
\]
is nonnegative, normalized, unbiased and pairwise independent. Its
total variation distance from $P$ is at most $B/(1+B)$, and
\begin{equation}\label{collision}
 C(Q)=\frac1M+
 \frac{C(P)-(1+\sum_{1\le|S|\le2}m_S^2)/M}{(1+B)^2}
 \ \ge\ \frac{C(P)}{(1+B)^2}.
\end{equation}
\end{theorem}
\begin{proof}
Nonnegativity is immediate from $B\ge f$. Uniform cube orthogonality
gives $\sum f=0$ and $M^{-1}\sum f\chi_S=m_S$ for the specified
sets $S$, so the low-order moments cancel exactly. If $B>0$, define
$R(x)=(B-f(x))/(MB)$, a probability law. Then
$Q=(P+BR)/(1+B)$, proving the variation bound and pointwise
$Q\ge P/(1+B)$, hence the collision lower bound. If $B=0$, then
$f=0$ and $Q=P$. For the exact identity expand the square, using
$\sum P f=\sum m_S^2$ and $M^{-1}\sum f^2=\sum m_S^2$.
\end{proof}

The smallest nonnegative buffer in this same affine ansatz is exactly
\[
 \beta_* =\max\{0,\max_x(f(x)-MP(x))\}.
\]
Replacing $B$ by any $\beta\ge\beta_*$ preserves normalization,
exact low-order moment cancellation, and the equality in
\eqref{collision}. The mixture interpretation and its simple
variation/collision lower bounds require $\beta\ge B$; they must not
be asserted for the potentially smaller $\beta_*$. Thus the explicit
universal construction is conservative, not an optimal repair theorem.

\begin{theorem}[A sufficient input window for exact extraction]
Let $d=n-1$, $n\equiv1\pmod4$, $n\ge5$, and $N=n+3$.
If $C(P)\ge1/N$ and $B\le1/(2048N^3)$, the corrected law satisfies
\[
 0\le1/N-C(Q)\le1/(1024N^4).
\]
Consequently the quantitative residue-one extraction theorem of round6
constructs a Hadamard matrix of order $N$ from $Q$.
It suffices to replace the assumption on $B$ by the stronger,
easier-to-check condition $L\le1/(2048N^3)$.
\end{theorem}
\begin{proof}
Exact isotropy of the augmented rows $(1,x)$ follows from the
corrected moments, and the residue collision bound gives $C(Q)\le1/N$.
The first theorem gives
\[
 1/N-C(Q)\le\frac1N\left(1-\frac1{(1+B)^2}\right)
 \le 2B/N\le1/(1024N^4).
\]
Apply the round6 stability theorem. This implication assumes the two
stated input properties and does not produce such an input in unknown
Hadamard orders.
\end{proof}

\begin{theorem}[Quantization obstruction for a uniform row candidate]
In the same parameters, suppose $P$ is uniform on $N$ distinct sign
rows, and at least one first or second moment is nonzero. The universal
$B$-buffer correction above satisfies
\[
 1/N-C(Q)\ge1/(2N^2),
\]
so it cannot enter the round6 extraction window.
\end{theorem}
\begin{proof}
Since $N$ is even, each moment is an even integer divided by $N$;
a nonzero one has magnitude at least $2/N$. Uniform expectation
gives $|m_S|=|\mathbb E_{\rm unif} f\chi_S|
\le\mathbb E_{\rm unif}|f|\le2B$, since $f$ has mean zero
and positive part at most $B$. Thus $B\ge1/N$.
Here $C(P)=1/N$ and $M=2^{N-4}\ge2N$ for $N\ge8$.
The exact collision identity yields
\[
 1/N-C(Q)\ge(1/N-1/M)\left(1-(1+B)^{-2}\right)
 \ge\frac1{2N}\left(1-(1+1/N)^{-2}\right)
 \ge\frac1{2N^2}.
\]
The last inequality is equivalent to $N^2-N-1\ge0$.
\end{proof}
This obstruction concerns the specified guaranteed buffer, not every
correction scheme or the smaller $\beta_*$. It explains why repairing
an arbitrary nearly orthogonal nearest-order sign matrix by adding a
uniform Fourier buffer is inadequate. For uniform $N$-row candidates,
the sufficient error condition of the preceding theorem forces all
low-order moments to vanish already by integer quantization.
\end{document}
